% ---- pg-583 （书页 571） ----

特斯一致性残差平方和量度, 以及如何确定 $\mathbf{Y}^{*}$ 相对 $\mathbf{X}^{*}$ 的普罗克鲁斯特斯旋转.

\mvthm{结论 12.2}{设 $n\times p$ 结构 $\mathbf{X}^{*}$ 和 $\mathbf{Y}^{*}$ 均已中心化, 即各行均值皆为零, 此时}
\phantomsection\label{eq:12-50}
\begin{equation}
\begin{aligned}
PR^{2}&=\sum_{j=1}^{n}\mathbf{x}_{j}'\mathbf{x}_{j}+\sum_{j=1}^{n}\mathbf{y}_{j}'\mathbf{y}_{j}-2\sum_{i=1}^{p}\lambda_{i}\\
&=\mathrm{tr}[\mathbf{X}^{*}\mathbf{X}^{*\prime}]+\mathrm{tr}[\mathbf{Y}^{*}\mathbf{Y}^{*\prime}]-2\mathrm{tr}[\boldsymbol{\Lambda}]
\end{aligned}
\tag{12-50}
\end{equation}
其中 $\boldsymbol{\Lambda}=$ 对角$(\lambda_{1},\lambda_{2},\cdots,\lambda_{p})$, 最小化变换为
\phantomsection\label{eq:12-51}
\begin{equation}
\hat{\mathbf{Q}}=\sum_{i=1}^{p}\mathbf{v}_{i}\mathbf{u}_{i}'=\mathbf{V}\mathbf{U}'\qquad \hat{\mathbf{b}}=\mathbf{0} \tag{12-51}
\end{equation}
这里 $\boldsymbol{\Lambda},\mathbf{U}$ 和 $\mathbf{V}$ 由下面的奇异值分解得到:
\[
\sum_{j=1}^{n}\mathbf{y}_{j}\mathbf{x}_{j}'=\underset{(p\times n)}{\mathbf{Y}^{*\prime}}\underset{(n\times p)}{\mathbf{X}^{*}}=\underset{(p\times p)}{\mathbf{U}}\underset{(p\times p)}{\boldsymbol{\Lambda}}\underset{(p\times p)}{\mathbf{V}'}
\]

证明\quad 由于两个点结构是已实现零均值中心化$\left(\sum_{j=1}^{n}\mathbf{x}_{j}=\mathbf{0}, \text{及} \sum_{j=1}^{n}\mathbf{y}_{j}=\mathbf{0}\right)$, 我们有
\[
\sum_{j=1}^{n}(\mathbf{x}_{j}-\mathbf{Q}\mathbf{y}_{j}-\mathbf{b})'(\mathbf{x}_{j}-\mathbf{Q}\mathbf{y}_{j}-\mathbf{b})=\sum_{j=1}^{n}(\mathbf{x}_{j}-\mathbf{Q}\mathbf{y}_{j})'(\mathbf{x}_{j}-\mathbf{Q}\mathbf{y}_{j})+n\mathbf{b}'\mathbf{b}
\]
右端的最后一项为非负项, 因此当 $\hat{\mathbf{b}}=\mathbf{0}$ 有最佳拟合. 因此我们只需考虑
\[
PR^{2}=\min_{\mathbf{Q}}\sum_{j=1}^{n}(\mathbf{x}_{j}-\mathbf{Q}\mathbf{y}_{j})'(\mathbf{x}_{j}-\mathbf{Q}\mathbf{y}_{j})=\sum_{j=1}^{n}\mathbf{x}_{j}'\mathbf{x}_{j}+\sum_{j=1}^{n}\mathbf{y}_{j}'\mathbf{y}_{j}-2\max_{\mathbf{Q}}\sum_{j=1}^{n}\mathbf{x}_{j}'\mathbf{Q}\mathbf{y}_{j}
\]
利用 $\mathbf{x}_{j}'\mathbf{Q}\mathbf{y}_{j}=\mathrm{tr}[\mathbf{Q}\mathbf{y}_{j}\mathbf{x}_{j}']$, 可将上式右端最大化的表达式变成
\[
\sum_{j=1}^{n}\mathbf{x}_{j}'\mathbf{Q}\mathbf{y}_{j}=\sum_{j=1}^{n}\mathrm{tr}[\mathbf{Q}\mathbf{y}_{j}\mathbf{x}_{j}']=\mathrm{tr}\left[\mathbf{Q}\sum_{j=1}^{n}\mathbf{y}_{j}\mathbf{x}_{j}'\right]
\]
由奇异值分解,
\[
\sum_{j=1}^{n}\mathbf{y}_{j}\mathbf{x}_{j}'=\mathbf{Y}^{*\prime}\mathbf{X}^{*}=\mathbf{U}\boldsymbol{\Lambda}\mathbf{V}'=\sum_{i=1}^{p}\lambda_{i}\mathbf{u}_{i}\mathbf{v}_{i}'
\]
其中 $\mathbf{U}=[\mathbf{u}_{1},\mathbf{u}_{2},\cdots,\mathbf{u}_{p}],\mathbf{V}=[\mathbf{v}_{1},\mathbf{v}_{2},\cdots,\mathbf{v}_{p}]$ 为 $p\times p$ 正交矩阵. 因此,
\[
\sum_{j=1}^{n}\mathbf{x}_{j}'\mathbf{Q}\mathbf{y}_{j}=\mathrm{tr}\left[\mathbf{Q}\left(\sum_{i=1}^{p}\lambda_{i}\mathbf{u}_{i}\mathbf{v}_{i}'\right)\right]=\sum_{i=1}^{p}\lambda_{i}\mathrm{tr}[\mathbf{Q}\mathbf{u}_{i}\mathbf{v}_{i}']
\]
上式右端第 $i$ 项中可变的量
\[
\mathrm{tr}[\mathbf{Q}\mathbf{u}_{i}\mathbf{v}_{i}']=\mathbf{v}_{i}'\mathbf{Q}\mathbf{u}_{i}
\]
有上界 1. 这可以用柯西-施瓦茨不等式[见\hyperref[eq:2-48]{式(2-48)}]并令 $\mathbf{b}=\mathbf{Q}\mathbf{v}_{i}$ 和 $\mathbf{d}=\mathbf{u}_{i}$ 来证明. 也就是说
\[
\mathbf{v}_{i}'\mathbf{Q}\mathbf{u}_{i}\leqslant\sqrt{\mathbf{v}_{i}'\mathbf{Q}'\mathbf{Q}\mathbf{v}_{i}}\sqrt{\mathbf{u}_{i}'\mathbf{u}_{i}}=\sqrt{\mathbf{v}_{i}'\mathbf{v}_{i}}\times 1=1
\]
所有 $p$ 个这样的项, 都可以令 $\hat{\mathbf{Q}}=\mathbf{V}\mathbf{U}'$ 使之达到最大. 在这种选择下,
\[
\mathbf{v}_{i}'\hat{\mathbf{Q}}\mathbf{u}_{i}=\mathbf{v}_{i}'\mathbf{V}\mathbf{U}'\mathbf{u}_{i}
=[0,\cdots,0,1,0,\cdots,0]
\begin{bmatrix}
0\\ \vdots\\ 0\\ 1\\ 0\\ \vdots\\ 0
\end{bmatrix}=1
\]
因此
